\section{Formal verification}\label{sec:formalization}

This section records the correspondence between the paper statements and
their formal evidence; it is not a second proof narrative.  The structural
reductions are ordinary Lean proofs.  A separate Mathlib-only package checks
the exceptional two-fixed-point Q25 census in the normalized coordinate model
used inside the proof of Proposition~\ref{prop:f2}; the projective
normalizations and semantic transport are manuscript arguments.  The
development uses Lean~\cite{Lean2015} on top of mathlib~\cite{Mathlib2020}
and is divided as follows.

\begin{enumerate}[leftmargin=2.2em]
\item The projective-conjugation and quadratic-Frobenius modules construct
      the semilinear projective involution, prove the degree-two fixed-locus
      normalization, and count the $N=(s^2-s)/2$ candidates on a fixed line.
\item The line-counting and forbidden-candidate modules prove
      Lemma~\ref{lem:empty-lines}, the exact secant-orbit count, the
      injective charge, and the semantic arc-extension conclusion in
      Theorem~\ref{thm:pair-criterion}.
\item The global-count module defines the finite set of fresh nonfixed
      Frobenius pairs whose union with the old arc remains an arc.  It proves
      that this semantic set is the disjoint union of the carrierwise legal
      sets and that its cardinality is the abstract legal count.
\item The alternate-repair module formalizes deletion of an arbitrary
      selected nonfixed orbit, proves that the erased orbit is one semantic
      legal pair of the invariant eight-arc remainder, and supplies the
      certificate-free eight-repair baseline.  Its order-five wrapper proves
      Corollary~\ref{cor:q25-nonexceptional-repair}.
\item The profile-envelope modules prove the five closed forms, their
      piecewise lower envelope, the semantic envelope comparison in
      Theorem~\ref{thm:profile-envelope}, and the uniform $318$-repair
      conclusion in Corollary~\ref{cor:alternate-repair}.
\item The phase-diagram and parameterized-repair modules prove the exact
      maximum $W(k)=\lfloor(k-1)^2/4\rfloor$, derive the empty-carrier range
      from the phase inequality itself, and prove
      Theorem~\ref{thm:parameterized-repair} for semantic legal-pair finsets.
\item The collision and invisibility modules prove
      Proposition~\ref{prop:linewise}, Corollary~\ref{cor:aggregate}, the
      fixed-center interpretation, and the cross-pair estimate
      \eqref{eq:cross-invisible}.
\item Separate profile modules prove the lower bounds $5$ and $4$ for
      $f=0$ and $f=4$.  For $f=2$, residual composition certificates prove
      the lower bound $32$ on all $1189$ representatives; orbit--stabilizer,
      strict-composition, and exhaustive-dispatch modules prove attainment
      and equality-orbit exhaustion in the package's normalized coordinate
      model.  The proof of Proposition~\ref{prop:f2} carries these conclusions
      through both normalizations.  The order-five
      alternate-repair module and the paper-level parity split then give
      Theorem~\ref{thm:q25} and
      Corollary~\ref{cor:q25-uniform-repair}.
\item The orbit-saturation module proves the denominator-free inequality
      $2s(s-1)\le(k-1)^2$ used in Corollary~\ref{cor:saturation}.  The final
      ceiling notation and the geometric case split are paper proofs.
\end{enumerate}

Corollary~\ref{cor:clebsch-pair-extension} is a manuscript-level
specialization of the semantic carrier count to $(s,f,e)=(11,6,0)$.  Its
arithmetic and the observation $M=0$ are printed in full; no separate Lean
terminal is claimed for the number $4180$.

The scoped structural builds and the sealed certificate package's
declaration-level axiom audits report exactly
\[
  [\texttt{propext},\ \texttt{Classical.choice},\ \texttt{Quot.sound}],
\]
the standard classical quotient profile inherited from mathlib.  No public
theorem depends on an admitted proof, a custom axiom, or
\texttt{native\_decide}.  The finite leaves supporting
Proposition~\ref{prop:f2} use kernel \texttt{decide}.  The certificate terminal
is the normalized classification; the passage to semantic arcs is the
manuscript proof and does not rely on a count of semantic arcs.

The certificate aggregate exposes
\path{PairCertificate.allRows},
\path{ResidualMinimumOrbits.card_minimumOrbitUnion}, and the two normalized
exhaustion terminals in \path{Exhaustion.lean}.

For reproducibility, the paper-facing declarations are located as follows.

\begin{center}
\small
\begin{tabular}{>{\raggedright\arraybackslash}p{0.43\linewidth}
                >{\raggedright\arraybackslash}p{0.49\linewidth}}
\toprule
Source module & Paper-facing declaration or role \\
\midrule
\path{FiniteGeom/BaerCompletion/PairExtension.lean} &
  \path{quadraticBaer_pairExtension_lowerBound} \\
\path{RelativeConicArcs/QuadraticForbidden.lean} &
  \path{exists_quadratic_pair_extension} \\
\path{RelativeConicArcs/QuadraticGlobalCount.lean} &
  \path{card_globalLegalPairs_eq_legalCount} \\
\path{RelativeConicArcs/AlternateOrbitRepair.lean} &
  \path{eight_le_alternateLegalPairs_of_seven_le} \\
\path{RelativeConicArcs/AlternateOrbitRepairProfileEnvelope.lean} &
  five closed forms and the piecewise profile envelope \\
\path{RelativeConicArcs/}\newline
\path{AlternateOrbitRepairProfile}\allowbreak\path{EnvelopeResult.lean} &
  \path{profileEnvelope_le_card_globalLegalPairs_of_card_eight},
  \path{three_hundred_eighteen_le_alternateLegalPairs_of_seven_le} \\
\path{RelativeConicArcs/}\newline
\path{AlternateOrbitRepairPhaseDiagram.lean} &
  \path{exists_profile_attaining_worstObstruction},
  \path{phase_profile_product_lowerBound} \\
\path{RelativeConicArcs/Parameterized}\newline
\path{AlternateOrbitRepair.lean} &
  \path{card_alternateLegalPairs_ge_of_phase} \\
\path{RelativeConicArcs/QuadraticInvisible.lean} &
  aggregate equality/excess and fixed-center criteria \\
\path{FiniteGeom/BaerCompletion/OrbitSaturation.lean} &
  \path{orbitSaturation_quadratic_bound_of_split} \\
\bottomrule
\end{tabular}
\end{center}

\subsection*{Artifact and reproduction}

The exceptional finite layer uses a separately versioned Mathlib-only
certificate package.  Its exact commit and manifest digest are recorded in
the verification pin.
That package pins Lean 4.32.0-rc1 and mathlib revision
\texttt{571b8a8e5421}.  From the package root
its sealed aggregate and source-reproduction checks are run by
\begin{verbatim}
nix run .#verify
\end{verbatim}
The exact pin and trust boundary are recorded in
\path{verification/q25-certificate-pin.json} and
\path{verification/README.md}.  The standalone paper archive does not vendor
the shared human-scale Lean development listed above; its structural theorem
sources and the separately pinned finite certificate have distinct release
boundaries.
The manuscript PDF is reproduced from its source directory with
\begin{verbatim}
nix shell nixpkgs#tectonic -c tectonic \
  equivariant-robust-completion.tex --keep-logs
\end{verbatim}

\clearpage
